\metatitle{Lattice path matroids}
\metadescription{Introduction to lattice path matroids, including their relation to transversal matroids, Tutte polynomials, skew Ferrers shapes, Dyck paths, and width-two posets.}

\section[latticePathMatroids]{Lattice path matroids}


The family of \defin{lattice path matroids} was introduced by
\name{Joseph Bonin}, \name[A. de Mier]{Anna de Mier} and \name{Marc Noy}
\cite{BoninMierNoy2003}. This family is closed under all \hyperref[matroidMinors]{matroid minors}
(deletions and contractions) and also under \hyperref[matroidDuality]{matroid duality}.
All lattice path matroids are \hyperref[transversalMatroids]{transversal matroids}.
They are also positroids, see \cite{Oh2011}.

\name{Yiming Chen}, \name{Yao Li}, and \name{Ming Yao} study the
$(a,b)$-Catalan subfamily through valuative invariants
\cite{ChenLiYao2025x}.  They express the indicator function of an
$(a,b)$-Catalan matroid polytope as a weighted sum of indicator functions of
direct sums of uniform matroids.  This gives explicit formulas for valuative
invariants, including Ehrhart positivity and Kazhdan--Lusztig invariants for
these Catalan matroids.

The \hyperref[tuttePolynomial]{Tutte polynomial} for a lattice path matroid
can be computed efficiently, see \cite{BoninMierNoy2003}.
Moreover, there is a nice combinatorial description of
the internal and external activity for the bases.


\begin{example}
\begin{figure}
 \svgimg[width=0.2\textwidth]{./svg-images/latticepath-matroid.svg}{A skew Ferrers shape and labeling of the East steps.}
\end{figure}
\end{example}


For an introduction to lattice path matroids,
see \href{https://maa.org/sites/default/files/pdf/shortcourse/2011/TransversalNotes.pdf}{these lecture notes}
by \name{Joseph Bonin}.


\begin{example}[Lattice path matroid from a Dyck path]
Consider the interval family
\[
A_1 = \{1,2,3,4,5,6\},\qquad
A_2 = \{3,4,5,6,7,8\},\qquad
A_3 = \{5,6,7,8,9,10\},\qquad
A_4 = \{6,7,8,9,10,11\},\qquad
A_5 = \{10,11,12\}.
\]
The corresponding skew shape is the Young diagram $\lambda/\mu$ with
\(
\lambda = (7,7,7,7,6),\qquad \mu = (5,2,2,1).
\)
Numbering the skew boxes by the elements of $A_5,\dots,A_1$ gives
\begin{figure}
\svgimg[width=0.2\textwidth]{./svg-images/latticepath-matroid2.svg}{The skew Ferrers shape and labeling of the North steps.}
\end{figure}


The upper boundary has North-step set $\{1,3,5,6,10\}$ and East-step set $\{2,4,7,8,9,11,12\}$.

The corresponding naturally labeled width-two permutation poset has chains
\[
  C_N = \{1 \lt 3 \lt 5 \lt 6 \lt 10\},\qquad
  C_E = \{2 \lt 4 \lt 7 \lt 8 \lt 9 \lt 11 \lt 12\},
\]
together with the cover relations $2\lt 3$, $4 \lt 5$, $9 \lt 10$ and $1 \lt 12$.
\begin{figure}
 \svgimg{./svg-images/latticepath-width-two-poset.svg}{The width-two poset associated with the lattice path matroid example.}
\end{figure}
The permutation $\pi = [2,4,7,8,9,11,1,12,3,5,6,10]$ defines a permutation poset,
and all linear extensions of the width-2-poset above are precisely the permutations below $\pi$
in the \hyperref[weakOrder]{weak order}. Note that $\pi$ is 321-avoiding.
Thus, for an ordinary Ferrers shape, the determinant in
\hyperref[LGVLatticePathsInSkewShape]{the LGV section on lattice paths in skew shapes}
computes the number of linear extensions of the associated width-two permutation poset.


We can also model set systems using a Dyck path.
Each set in the set system corresponds to a peak.
The peaks are marked with the corresponding set and the shaded boxes are those below and to the right of each peak.
\begin{figure}
\begin{ytableau}
 &  &  &  &  &  &  &  &  & *(lightgray)A_5 & *(lightgray) & *(lightblue)12 \\
 &  &  &  &  & *(lightgray)A_4 & *(lightgray) & *(lightgray) & *(lightgray) & *(lightgray) & *(lightblue)11 \\
 &  &  &  & *(lightgray)A_3 & *(lightgray) & *(lightgray) & *(lightgray) & *(lightgray) & *(lightblue)10 \\
 &  &  &  & *(lightgray) & *(lightgray) & *(lightgray) & *(lightgray) & *(lightblue)9 \\
 &  & *(lightgray)A_2 & *(lightgray) & *(lightgray) & *(lightgray) & *(lightgray) & *(lightblue)8 \\
 &  & *(lightgray) & *(lightgray) & *(lightgray) & *(lightgray) & *(lightblue)7 \\
*(lightgray)A_1 & *(lightgray) & *(lightgray) & *(lightgray) & *(lightgray) & *(lightblue)6 \\
*(lightgray) & *(lightgray) & *(lightgray) & *(lightgray) & *(lightblue)5 \\
*(lightgray) & *(lightgray) & *(lightgray) & *(lightblue)4 \\
*(lightgray) & *(lightgray) & *(lightblue)3 \\
*(lightgray) & *(lightblue)2 \\
*(lightblue)1
\end{ytableau}
\end{figure}
A basis of the lattice path matroid is then a selection of squares on the diagonal, $d_1,d_2,\dotsc,d_k$
such that peak $i$ "can see" the square $d_i$. Here, $k$ is the number of peaks.


In this model, the rank of the lattice path matroid is the number of peaks of the Dyck path,
and the dimension of its base polytope is the number of non-zero entries in the area sequence.
Area sequences of the form $[0,1,2,\dotsc,k,k,\dotsc,k]$ correspond to uniform matroids.
A combinatorial formula for the $h^*$-vector is provided in \cite{Kim2020}, see also \cite{Li2012}.


%If the Dyck path contains $N^2 E N^2 E^2 N E^2$ as a subsequence, then we have Q6.
\end{example}

The corresponding determinant formula for the lattice-path-matroid basis generating
function is given in \hyperref[LGVLatticePathsInSkewShape]{the LGV section on lattice paths in skew shapes}.


\section[rookMatroids]{Rook matroids}

Rook matroids, introduced in \cite{AlexanderssonJal2024x}, give another
matroidal structure on the same skew Ferrers shapes. They form a family of
\hyperref[transversalMatroids]{transversal matroids} that is compatible with
\hyperref[matroidDuality]{duality}, but whose behavior under
\hyperref[matroidMinors]{minors} is subtler than for lattice path matroids.
They also have the same \hyperref[tuttePolynomial]{Tutte polynomial} as the
corresponding lattice path matroid.

Let $\lambda/\mu$ be a skew Ferrers board with $r$ rows and $c$ columns.
Label the rows by $1,\dotsc,r$ from top to bottom, and label the columns by
$r+1,\dotsc,r+c$ from left to right. A rook placement is
\defin{non-attacking} if no two rooks share a row or a column. Two rooks form a
\defin{nesting} if one lies strictly southeast of the other. A
\defin{non-nesting rook placement} is a non-attacking placement with no nesting.
Let
\[
  \mathrm{NN}_{\lambda/\mu}
\]
denote the set of non-nesting rook placements on $\lambda/\mu$.
The corresponding \defin{non-nesting rook polynomial} is
\[
  M_{\lambda/\mu}(t) = \sum_k r_k(\lambda/\mu)t^k,
\]
where $r_k(\lambda/\mu)$ is the number of non-nesting rook placements with $k$
rooks.

For a placement $\rho \in \mathrm{NN}_{\lambda/\mu}$, let $R(\rho)$ be the set
of unoccupied row labels and let $C(\rho)$ be the set of occupied column labels.
The \defin{rook matroid} of $\lambda/\mu$ has bases
\[
  \mathcal{R}_{\lambda/\mu}
  =
  \{ R(\rho) \cup C(\rho) : \rho \in \mathrm{NN}_{\lambda/\mu} \}.
\]
Every such set has size $r$.

\begin{theorem}[Alexandersson--Jal, \cite{AlexanderssonJal2024x}]
For every skew shape $\lambda/\mu$, the collection
$\mathcal{R}_{\lambda/\mu}$ is the set of bases of a
\hyperref[transversalMatroids]{transversal matroid}.
\end{theorem}

\begin{example}
For the skew shape $\lambda/\mu = (4,3,2)/(1)$, the rows are labeled
$1,2,3$ and the columns are labeled $4,5,6,7$. The placement below is
non-nesting:
\begin{ytableau}
\none & \none 4 & \none 5 & \none 6 & \none 7 \\
\none 1 & \none &  &  & \bullet \\
\none 2 &  &  &  \\
\none 3 & \bullet &
\end{ytableau}

Here row $2$ is unoccupied and columns $4$ and $7$ are occupied, so this
placement gives the basis $R(\rho)\cup C(\rho) = \{2,4,7\}.$
\end{example}

\begin{example}
With the row-labeling convention above, the rook matroid of a
$k\times(n-k)$ rectangle is the uniform matroid $U_{k,n}$. Equivalently, an
$(n-k)\times k$ rectangle gives $U_{n-k,n}$.
\end{example}

\begin{theorem}[Alexandersson--Jal, \cite{AlexanderssonJal2024x}]
Rook matroids are closed under direct sums and under \hyperref[matroidDuality]{matroid duality}. More
precisely,
\[
  \mathcal{R}_{\lambda/\mu}^{*}\cong
  \mathcal{R}_{\lambda'/\mu'}.
\]
However, the class of rook matroids is not closed under all minors.
\end{theorem}

\begin{theorem}[Alexandersson--Jal, \cite{AlexanderssonJal2024x}]
Every rook matroid is a \hyperref[positroid]{positroid}.
\end{theorem}
Thus both lattice path matroids and rook matroids belong to the class of
positroids, even though the two matroid structures on a skew shape are not
always isomorphic.



\subsection[rookLatticePathMatroids]{Relation with lattice path matroids}

Write $\mathcal{P}_{\lambda/\mu}$ for the lattice path matroid whose upper and
lower boundary paths bound the skew shape $\lambda/\mu$. There is always a
bijection between lattice paths in $\lambda/\mu$ and non-nesting rook
placements in $\lambda/\mu$: place a rook in each valley of the path that lies
strictly inside the shape. The bijection is not always a matroid isomorphism.

\begin{theorem}[Alexandersson--Jal, \cite{AlexanderssonJal2024x}]
The rook matroid $\mathcal{R}_{\lambda/\mu}$ is isomorphic to the lattice path
matroid $\mathcal{P}_{\lambda/\mu}$ if and only if the skew shape
$\lambda/\mu$ avoids the subshape $332/1$.
\end{theorem}

Equivalently, the obstruction is the matroid $Q_6$: one has
$\mathcal{R}_{332/1}\cong Q_6$, and $Q_6$ is an excluded minor for lattice path
matroids \cite{Bonin2010}. Thus the rook and lattice-path structures agree for
many shapes, but not for all skew shapes.

Under the skew-shape--poset correspondence described below, the same condition
can be read as avoidance of an induced cycle poset. Namely, the associated
width-two poset has no induced subposet on two three-element chains
\[
  a_1 \lt a_2 \lt a_3,\qquad b_1 \lt b_2 \lt b_3
\]
whose only cross-chain cover relations are
\[
  a_1 \lt b_3,\qquad b_1 \lt a_3.
\]
Equivalently, the forbidden induced poset is the subdivision of the
four-element crown $C_2$ obtained by inserting one element on two opposite
covers. The two crossing cover relations correspond to the outer and inner
corners of the $332/1$ subshape.

In the other direction, every lattice path matroid is obtained as a deletion of
a sufficiently enlarged rook matroid. If $\lambda$ has $r$ rows, define
\[
  \lambda^\ast =
  (\lambda_1+r,\lambda_2+r-1,\dotsc,\lambda_r+1),
\]
and similarly add the staircase $(r-1,r-2,\dotsc)$ to $\mu$ to obtain
$\mu^\ast$. Then
\[
  \mathcal{P}_{\lambda/\mu}
  \cong
  \mathcal{R}_{\lambda^\ast/\mu^\ast}\setminus [r].
\]

Even when the two matroids are not isomorphic, their Tutte polynomials agree:
\[
  T(\mathcal{R}_{\lambda/\mu};x,y)
  =
  T(\mathcal{P}_{\lambda/\mu};x,y).
\]

\subsection[rookLatticePathPolytopes]{Base polytopes and valuative invariants}

For any matroid $M$, the \hyperref[matroidPolytopes]{matroid base polytope}
is the convex hull of the indicator vectors of its bases. Thus the rook
matroid base polytope is
\[
  \operatorname{conv}
  \{ \mathbf{e}_{R(\rho)\cup C(\rho)} :
     \rho\in \mathrm{NN}_{\lambda/\mu} \}.
\]
The base polytopes of lattice path matroids have an alcoved description and
snake decompositions, see \cite{BenedettiKnauerValencia2023x}.

The rook polytope also has a description in row-column coordinates. Besides
the inequalities $0\leq x_i,y_j\leq 1$ and the rank equation
\[
  x_1+\dotsb+x_r+y_1+\dotsb+y_c = r,
\]
the nontrivial inequalities are indexed by the inner and outer corners of the
skew shape. If the first $i$ rows occupy the consecutive columns
$a_i,\dotsc,b_i$, an inner corner gives an inequality of the form
\[
  (x_1+\dotsb+x_i) + (y_{a_i}+\dotsb+y_{b_i}) \leq b_i-a_i+1.
\]
Similarly, if rows $j,j+1,\dotsc,r$ occupy the consecutive columns
$a'_j,\dotsc,b'_j$, an outer corner gives
\[
  (x_j+\dotsb+x_r) + (y_{a'_j}+\dotsb+y_{b'_j}) \leq b'_j-a'_j+1.
\]

Bonin and de Mier proved a stronger relation between the two matroids:
$\mathcal{P}_{\lambda/\mu}$ and $\mathcal{R}_{\lambda/\mu}$ have the same
configuration \cite{BoninMier2026}. Consequently they have the same
Derksen--Fink $\mathcal{G}$-invariant, and every valuative matroid invariant
takes the same value on them. This implies, in particular, that the Tutte
polynomials agree and that the \hyperref[ehrhart]{Ehrhart polynomials} of the two matroid base
polytopes agree, even when the matroids are not isomorphic.

The base polytopes of lattice path matroids have positive Ehrhart coefficients, see \cite{FerroniMoralesPanova2026x}.




\subsection[rookMatroidPolynomials]{Rook polynomials and width-two posets}

The basis-generating polynomial of a rook matroid is
\hyperref[lorentzianPolynomial]{Lorentzian}. Consequently, the coefficient
sequence of $M_{\lambda/\mu}(t)$ is ultra-log-concave with no internal zeros
\cite{BrandenHuh2020,AlexanderssonJal2024x}. Real-rootedness can still fail.

The same paper gives a stability application: there is a generalized Catalan
matroid, hence a lattice path matroid, whose basis-generating polynomial is not
\hyperref[stablePolynomial]{stable}. Equivalently, this lattice path matroid
does not have the half-plane property, answering a natural question about
stability for this family of transversal matroids
\cite{ChoeOxleySokalWagner2004,ChoeWagner2006}.

The poset here is the same naturally labeled width-two poset associated with
the lattice paths above. Under the bijection between lattice paths and
non-nesting rook placements, the rook placements correspond to linear
extensions of this poset $P$, and
\[
  M_{\lambda/\mu}(t) = W_P(t),
\]
where $W_P(t)$ is the $P$-Eulerian polynomial. This refines the width-two
poset example above and is compatible with the \hyperref[weakOrder]{weak order}
description for ordinary Ferrers shapes.
